% Book 3, Year 1, chapter 1 -- the compilation pipeline (C17, gcc).
% Provenance (never printed): sources/book3/BRIEFS.md, chapter 1.
\chapter{From Source Code to a Running Program}\label{ch:b1:compilation-pipeline}

Type a short program into a file and run it with Python: the command \code{python hello_python.py} runs it
at once. Type the same program in C into \code{hello.c}, and nothing runs until another command,
\code{gcc hello.c -o hello}, has turned the text into a new file; only that file can be run. Between the text
you write and the program that runs there are four separate tools, three kinds of file and one layout of
memory. This chapter opens each of them, on this year's working language, C, and its \omterm{def:b1:compilation-pipeline:compiler}{compiler}, \code{gcc}.

\begin{recall}
A program is a text in a programming language, and the operating system starts it, gives it memory and
stops it. The processor only executes \emph{machine language}: numbered instructions stored in memory, as
the toy machine of the school volumes showed. An \emph{assembly language} names those instructions with
words and labels. You have also written Python: an \omterm{def:b1:compilation-pipeline:interpreter}{interpreter} reads the text and carries it out.
\end{recall}

\section{The four stages}\label{sec:b1:compilation-pipeline:stages}

Here is the C program we follow through every stage. It prints one line, using two \emph{\omterm{def:b1:compilation-pipeline:source-file}{macros}} written with
\texttt{\#define}.

\omcode{code/bachelor-1/01-compilation-pipeline/listings/hello.c}{2}{12}{A first C program}\label{lst:b1:compilation-pipeline:hello}

\omoutput{code/bachelor-1/01-compilation-pipeline/outputs/hello.txt}{One command builds the program, a second runs it.}

\begin{definition}[Source file, preprocessor]\label{def:b1:compilation-pipeline:source-file}\label{def:b1:compilation-pipeline:preprocessor}
A \emph{source file}\index{source file} is a text file written in a programming language, here a \code{.c} file
of C. The \emph{preprocessor}\index{preprocessor} is the first tool of a C build: it copies the source file,
replacing every \texttt{\#include} line by the text of the file it names, every use of a \emph{macro}\index{macro
(C)} (a name defined by \texttt{\#define}) by the text it stands for, and deleting comments. Its output is still C,
with no \texttt{\#} line left except line markers.
\end{definition}

\begin{definition}[Compiler, object file, linker, executable]\label{def:b1:compilation-pipeline:compiler}\label{def:b1:compilation-pipeline:object-file}\label{def:b1:compilation-pipeline:linker}\label{def:b1:compilation-pipeline:executable}
A \emph{compiler}\index{compiler} translates a program from one language into another, here from C into the
assembly language of the processor; the translation is called \emph{compilation}\index{compilation}. The
\emph{assembler} then encodes the assembly as machine code in an \emph{object file}\index{object file}: machine
code with holes, together with a table of the names it defines and the names it uses. Each name in that table
is a \emph{symbol}\index{symbol (of an object file)}. The \emph{linker}\index{linker} combines object files and
libraries into one \emph{executable file}\index{executable file}, a file the operating system can load and
start; this step is \emph{linking}\index{linking}, and a used symbol that no input defines is an
\emph{undefined reference}\index{undefined reference}.
\end{definition}

The command \code{gcc} runs all four tools in turn and deletes the intermediate files. Three options stop it
early and keep them: \code{-E} after preprocessing, \code{-S} after compiling, \code{-c} after assembling. The
session below takes \cref{lst:b1:compilation-pipeline:hello} through each stage by hand.

\omoutput{code/bachelor-1/01-compilation-pipeline/outputs/stages.txt}{The four stages, one at a
time.}\label{out:b1:compilation-pipeline:stages}

Read it from the top. The twelve lines of \code{hello.c} became 750 after preprocessing: \texttt{\#include
<stdio.h>} pulled in the declarations of the standard input--output library, and the line that interests us is
now line 748, with \code{GREETING} replaced by the string and \code{SQUARE(n)} by \code{((n) * (n))}. The
\omterm{def:b1:compilation-pipeline:compiler}{compiler}, asked for optimised code (\code{-O2}), then went further than the text: there is no multiplication in
the assembly, only an instruction \code{movl} that loads the constant 49 into a register, because \code{n} is always 7.
On this system \code{printf} also became \code{__printf_chk}, a checking variant that the system's \omterm{def:b1:compilation-pipeline:compiler}{compiler}
substitutes when it can see the format string; the first instruction, \code{endbr64}, marks a legal target for
indirect jumps. You will not write x86-64 assembly in this book (the Year 2 volume writes RISC-V assembly), but
reading the \omterm{def:b1:compilation-pipeline:compiler}{compiler}'s output is a skill: it tells you what your program really asks the machine to do.

The \omterm{def:b1:compilation-pipeline:compiler}{object file} \code{hello.o} is in the ELF format (Executable and Linkable Format), the format of \omterm{def:b1:compilation-pipeline:compiler}{object files}
and executables on Linux. Its \omterm{def:b1:compilation-pipeline:compiler}{symbol} table, printed by \code{nm}, shows a \code{T} for \code{main} (code,
\emph{text} in the old vocabulary, defined here), two \code{r} for the two strings (read-only data, local to
the file), and a \code{U} for \code{__printf_chk}: used, but defined somewhere else. Only after \omterm{def:b1:compilation-pipeline:compiler}{linking} with the
C library does \code{./hello} run.

\begin{omfigure}
\begin{tikzpicture}[
  file/.style={draw, rounded corners=2pt, fill=omCode!8, minimum width=16mm, minimum height=8mm,
               font=\small\ttfamily, align=center},
  tool/.style={font=\scriptsize\itshape},
  what/.style={font=\scriptsize, align=center, anchor=north},
  >={Stealth[length=2mm]}]
\foreach \n/\t/\x in {c/hello.c/0, i/hello.i/3.2, s/hello.s/6.4, o/hello.o/9.6, x/hello/12.8}
  \node[file] (\n) at (\x,0) {\t};
\draw[->] (c) -- node[below=1pt, font=\scriptsize\ttfamily] {-E} (i);
\draw[->] (i) -- node[below=1pt, font=\scriptsize\ttfamily] {-S} (s);
\draw[->] (s) -- node[below=1pt, font=\scriptsize\ttfamily] {-c} (o);
\draw[->] (o) -- (x);
\node[tool] at (1.6,0.65) {preprocessor};
\node[tool] at (4.8,0.65) {compiler};
\node[tool] at (8.0,0.65) {assembler};
\node[tool] at (11.2,0.65) {linker};
\node[what] at (0,-0.55) {C text};
\node[what] at (3.2,-0.55) {C text,\\expanded};
\node[what] at (6.4,-0.55) {assembly};
\node[what] at (9.6,-0.55) {machine code\\with holes};
\node[file, fill=white] (lib) at (8.0,-2.3) {C library};
\draw[->] (lib.east) -- (11.2,-2.3) -- (11.2,-0.08);
\node[draw, rounded corners=6pt, fill=omThm!10, minimum width=19mm, minimum height=8mm,
      font=\small, align=center] (p) at (12.8,-2.3) {running\\[-2pt]program};
\draw[->] (x) -- node[right, tool] {loaded} (p);
\end{tikzpicture}
\par\smallskip{\small The four stages of a C build. Each arrow is a separate program; \code{gcc} runs them in
turn, and the options \code{-E}, \code{-S} and \code{-c} stop it after the first, second or third. The
\omterm{def:b1:compilation-pipeline:compiler}{linker} fills the holes left in the \omterm{def:b1:compilation-pipeline:compiler}{object file} with the addresses of the library's code.}
\end{omfigure}

\begin{history}[C and Unix]
% ledger: b1:c-years, b1:c-bcpl, b1:c-unix-1973, b1:c17
C came into being between 1969 and 1973 at Bell Labs, alongside the first versions of the Unix operating system,
its most creative year being 1972. It descends from a language called B, itself derived from BCPL, and by early
1973 it was strong enough for the Unix kernel to be rewritten in C, for the PDP-11 minicomputer, during that
summer. An operating system written in a language that a \omterm{def:b1:compilation-pipeline:compiler}{compiler} could carry to other processors was new, and
it is why C and Unix spread together. The language is now an international standard; this book follows the
2018 edition, C17 (ISO/IEC 9899:2018), which the \code{-std=c17} option asks for.
\end{history}

\omimg{book3/pdp-11-40.jpg}{0.26\linewidth}{A PDP-11/40 minicomputer in the Technisches Museum Wien: the processor,
its front panel of switches and lights, and its tape drives in one cabinet. Photograph: Jeff Keyzer, CC BY-SA 2.0.}

\section{Object files, symbols and the linker}\label{sec:b1:compilation-pipeline:linker}

A program of a few thousand lines is never one file. It is split into \emph{modules}: a module is a pair of
files, a \emph{header} (\code{.h}) that declares what the module offers, and a \code{.c} file that defines it.
\cref{ch:b1:c-modules-make-tests} develops the method; here is the smallest example, a module that
squares numbers and counts how often it was asked.

\omcode{code/bachelor-1/01-compilation-pipeline/listings/square.h}{2}{8}{The module's header: what it offers}\label{lst:b1:compilation-pipeline:square-h}

\omcode{code/bachelor-1/01-compilation-pipeline/listings/square.c}{2}{10}{The module's implementation}\label{lst:b1:compilation-pipeline:square-c}

\omcode{code/bachelor-1/01-compilation-pipeline/listings/main.c}{2}{13}{A program that uses the module}\label{lst:b1:compilation-pipeline:main}

The header says that a function \code{square} exists and takes and returns a \code{long}; it says that an
\code{int} called \code{square_calls} exists somewhere (\code{extern}). That is all \code{main.c} needs to be
compiled on its own: the \omterm{def:b1:compilation-pipeline:compiler}{compiler} checks every call against the declaration and leaves a hole where the
address of \code{square} will go. Each \code{.c} file is compiled separately into its \omterm{def:b1:compilation-pipeline:compiler}{object file}, and the
\omterm{def:b1:compilation-pipeline:compiler}{linker} then matches every \code{U} of one file with a \code{T} (code) or \code{B}, \code{D} (data) of another.

\omoutput{code/bachelor-1/01-compilation-pipeline/outputs/linking.txt}{Separate \omterm{def:b1:compilation-pipeline:compiler}{compilation}, a successful link and
two failed ones.}\label{out:b1:compilation-pipeline:linking}

\begin{example}[Two link errors]\label{ex:b1:compilation-pipeline:link-errors}
In the session, \code{gcc main.o -o broken} forgets \code{square.o}: \code{main.o} uses \code{square} and
\code{square_calls} and nobody defines them, so the \omterm{def:b1:compilation-pipeline:compiler}{linker} reports two \omterm{def:b1:compilation-pipeline:compiler}{undefined references} (and two warnings
that follow from the same missing data \omterm{def:b1:compilation-pipeline:compiler}{symbol}). The last command links \code{twice.o}, which defines a second
\code{square}: now the \omterm{def:b1:compilation-pipeline:compiler}{symbol} is defined twice and the \omterm{def:b1:compilation-pipeline:compiler}{linker} cannot choose. Note that the \omterm{def:b1:compilation-pipeline:compiler}{compiler} accepted
every file: each error belongs to the link, where the files meet.
\end{example}

\begin{proposition}[When a link succeeds]\label{prop:b1:compilation-pipeline:link}
Consider \omterm{def:b1:compilation-pipeline:compiler}{object files} in which every \omterm{def:b1:compilation-pipeline:compiler}{symbol} is either defined (\code{T}, \code{D}, \code{B}, \code{R}) or used
(\code{U}), and ignore libraries. \omterm{def:b1:compilation-pipeline:compiler}{Linking} them succeeds if and only if every \omterm{def:b1:compilation-pipeline:compiler}{symbol} used by one of the files is
defined by exactly one of them, and no \omterm{def:b1:compilation-pipeline:compiler}{symbol} is defined by two of them.
\end{proposition}

\begin{proof}
The \omterm{def:b1:compilation-pipeline:compiler}{linker} builds one table from all the definitions. If a name is defined twice, it reports a multiple
definition; if a used name is in no definition, it reports an \omterm{def:b1:compilation-pipeline:compiler}{undefined reference}; these are its only two ways
to fail in this model. Conversely, when every used name has exactly one definition, each hole is filled with
the address of that definition, and the executable is complete.
\end{proof}

\begin{remark}[Local symbols and libraries]\label{rem:b1:compilation-pipeline:local}
A lower-case letter in \code{nm} (\code{t}, \code{b}, \code{r}) marks a \emph{local} \omterm{def:b1:compilation-pipeline:compiler}{symbol}: a function or
variable declared \code{static} at file level, invisible to the \omterm{def:b1:compilation-pipeline:compiler}{linker}'s matching, so two files may each have
their own. Libraries add one rule: the \omterm{def:b1:compilation-pipeline:compiler}{linker} takes from a library only the \omterm{def:b1:compilation-pipeline:compiler}{object files} that define
something still undefined. The C library is linked by default, which is why \code{printf} resolved silently.
A library can be copied into the executable (static \omterm{def:b1:compilation-pipeline:compiler}{linking}) or loaded when the program starts (a shared
library, the default for the C library on Linux).
\end{remark}

\begin{omfigure}
\begin{tikzpicture}[
  sym/.style={font=\footnotesize\ttfamily, anchor=west, inner sep=1pt},
  obj/.style={draw, rounded corners=2pt, fill=omCode!6, minimum width=40mm, minimum height=23mm},
  >={Stealth[length=2mm]}]
\node[obj] (m) at (0,0) {};
\node[font=\small\bfseries, anchor=south] at (m.north) {main.o};
\node[sym] (m1) at (-1.8,0.75) {T main};
\node[sym] (m3) at (-1.8,0.25) {U square};
\node[sym] (m4) at (-1.8,-0.25) {U square\_calls};
\node[sym] (m2) at (-1.8,-0.75) {U printf};
\node[obj] (s) at (6.5,0) {};
\node[font=\small\bfseries, anchor=south] at (s.north) {square.o};
\node[sym] (s1) at (4.7,0.25) {T square};
\node[sym] (s2) at (4.7,-0.25) {B square\_calls};
\node[obj, minimum height=10mm] (l) at (6.5,-3) {};
\node[font=\small\bfseries, anchor=north] at (l.south) {C library};
\node[sym] (l1) at (4.7,-3) {T printf};
\draw[->, omDef, thick] (m.east |- m3) -- (s1.west);
\draw[->, omDef, thick] (m.east |- m4) -- (s2.west);
\draw[->, omDef, thick] (m.east |- m2) to[out=0, in=180] (l1.west);
\end{tikzpicture}
\par\smallskip{\small The \omterm{def:b1:compilation-pipeline:compiler}{linker}'s work in the session above: each \omterm{def:b1:compilation-pipeline:compiler}{symbol} that \code{main.o} uses (\code{U}) is
matched with the one file that defines it. A \code{U} with no arrow is an \omterm{def:b1:compilation-pipeline:compiler}{undefined reference}; a name with two
definitions has two possible arrows, and the link fails.}
\end{omfigure}

\section{Warnings are a tool}\label{sec:b1:compilation-pipeline:warnings}

C lets many mistakes through: the language trusts the programmer. The \omterm{def:b1:compilation-pipeline:compiler}{compiler}, though, notices a great deal,
and says so when it is asked. Every program in this book is compiled with \code{-Wall -Wextra}, which turn on
most of its checks.

\begin{definition}[Compiler warning]\label{def:b1:compilation-pipeline:warning}
A \emph{compiler warning}\index{compiler warning} is a message about a construct that is legal but probably
wrong. Unlike an error, it does not stop the build; the option \code{-Werror} turns every warning into an error,
and the tests of this book are built that way.
\end{definition}

\omcode{code/bachelor-1/01-compilation-pipeline/listings/warnings.c}{2}{17}{Three bugs, all of them legal C}\label{lst:b1:compilation-pipeline:warnings}

\omoutput{code/bachelor-1/01-compilation-pipeline/outputs/warnings.txt}{The \omterm{def:b1:compilation-pipeline:compiler}{compiler} finds all three, and the
program still runs.}\label{out:b1:compilation-pipeline:warnings}

Each warning gives the file, the line and the column, quotes the line, and names the option that enabled it in
square brackets. The first says that \code{if (x = 0)} assigns 0 to \code{x} and tests the result, which is
always false; the comparison is \code{x == 0}. The second says that \texttt{\%d} expects an \code{int} while
\code{big} is a \code{long}: the behaviour is undefined, and here \code{printf} took the low 32 bits of the
\code{long} and printed $-1294967296$ instead of three billion. The third says that \code{count} is used before any value was stored in it. Its printed value, 3,
is luck: an uninitialised variable may hold anything, and the next \omterm{def:b1:compilation-pipeline:compiler}{compiler} version may print something else.
The fixed version (\code{warnings_fixed.c} in the chapter's code) compiles without a message and prints

\omoutput{code/bachelor-1/01-compilation-pipeline/outputs/warnings_fixed.txt}{The corrected program, built with no
warning.}

\begin{inpractice}[Read the first message first]
A single mistake can produce a cascade of messages. Fix the first one, rebuild, and read again; later messages
are often consequences. A build that prints nothing is the goal of every program in this book.
\end{inpractice}

\section{The memory of a running program}\label{sec:b1:compilation-pipeline:memory}

When the operating system starts an executable, it gives the running program its own memory, laid out in
regions. The program below creates one object of each kind, sorts their addresses and prints which comes first.
It never prints an address itself: the numbers change from machine to machine, the order does not.

\omcode{code/bachelor-1/01-compilation-pipeline/listings/regions.c}{2}{41}{Which object lies below which?}\label{lst:b1:compilation-pipeline:regions}

\omoutput{code/bachelor-1/01-compilation-pipeline/outputs/regions.txt}{The order of the regions on this
machine.}\label{out:b1:compilation-pipeline:regions}

\begin{definition}[Memory layout of a process]\label{def:b1:compilation-pipeline:layout}
The \emph{memory layout of a process}\index{memory layout of a process} is the arrangement of a running
program's memory in regions. The \emph{code segment}\index{code segment} holds the machine code and, next to it,
read-only data such as string literals; the \emph{data segment}\index{data segment} holds the variables that
exist for the whole run, the initialised ones first and then the zero-initialised ones (traditionally called
\emph{bss}). Above them grows the \emph{heap}, from which the program requests blocks while it runs
(\cref{ch:b1:dynamic-memory}); near the top of the address space is the \emph{call stack}, which holds the
local variables of the functions being executed (\cref{ch:b1:c-functions-and-stack}).
\end{definition}

\begin{omfigure}
\begin{tikzpicture}[
  seg/.style={draw, minimum width=46mm, align=center, font=\small},
  life/.style={font=\footnotesize, anchor=west, align=left},
  >={Stealth[length=2mm]}]
\node[seg, minimum height=9mm, fill=omCode!10] (code) at (0,0) {code and read-only data};
\node[seg, minimum height=8mm, fill=omDef!10, anchor=south] (data) at (code.north) {initialised data};
\node[seg, minimum height=8mm, fill=omDef!5, anchor=south] (bss) at (data.north) {zero-initialised data};
\node[seg, minimum height=11mm, fill=omProp!10, anchor=south] (heap) at (bss.north) {heap};
\node[seg, minimum height=16mm, anchor=south, draw=black!40, dashed] (gap) at (heap.north) {};
\node[seg, minimum height=11mm, fill=omThm!10, anchor=south] (stack) at (gap.north) {call stack};
\draw[->, thick] ([xshift=12mm]heap.north) -- ++(0,0.6) node[left, font=\scriptsize] {grows up};
\draw[->, thick] ([xshift=-12mm]stack.south) -- ++(0,-0.6) node[right, font=\scriptsize] {grows down};
\node[life] at (2.6,3.9) {unused};
\node[life] at (2.6,0) {whole run, read-only};
\node[life] at (2.6,0.85) {whole run};
\node[life] at (2.6,1.65) {whole run, starts at zero};
\node[life] at (2.6,2.6) {from \code{malloc} to \code{free}};
\node[life] at (2.6,5.3) {one function call};
\draw[->] (-3.2,-0.4) -- (-3.2,5.8) node[above, font=\scriptsize] {higher addresses};
\end{tikzpicture}
\par\smallskip{\small The regions of a running program, in the order \cref{out:b1:compilation-pipeline:regions}
found them, with how long an object in each region lives. The sizes are not to scale: the gap between the heap
and the stack is by far the largest part of the address space.}
\end{omfigure}

The lifetimes on the right of the figure are the real lesson. A global variable exists from the start of the run
to its end; a local variable exists only while its function runs, and its memory is reused by the next call; a
heap block lives exactly as long as the program decides. Most of the bugs of the next four chapters are a
pointer that outlives the object it points to.

\section{Interpreted and compiled}\label{sec:b1:compilation-pipeline:interpreted}

Here is the program of the opening, in Python. One command, and it runs.

\omcode{code/bachelor-1/01-compilation-pipeline/listings/hello_python.py}{2}{10}{The opening program in Python}

\omoutput{code/bachelor-1/01-compilation-pipeline/outputs/hello_python.txt}{No separate build step.}

Python does not skip the translation: it does it differently, and out of sight.

\begin{definition}[Interpreter, bytecode]\label{def:b1:compilation-pipeline:interpreter}
An \emph{interpreter}\index{interpreter} is a program that carries out another program by reading it, rather
than translating it in advance into machine code. CPython, the usual Python, first compiles each function into
\emph{bytecode}\index{bytecode}: instructions for an imaginary machine, which a loop written in C fetches one by
one and executes.
\end{definition}

\omcode{code/bachelor-1/01-compilation-pipeline/listings/bytecode.py}{2}{10}{Asking Python for its \omterm{def:b1:compilation-pipeline:interpreter}{bytecode}}

\omoutput{code/bachelor-1/01-compilation-pipeline/outputs/bytecode.txt}{The \omterm{def:b1:compilation-pipeline:interpreter}{bytecode} of
\code{square_plus}: \code{RESUME}, which starts the function, then six instructions.}\label{out:b1:compilation-pipeline:bytecode}

The function became a stack machine program: push \code{x}, push \code{x}, multiply, push \code{y}, add, return.
Each instruction is decoded by the \omterm{def:b1:compilation-pipeline:interpreter}{interpreter}'s loop at every execution, and the loop checks the types of
\code{x} and \code{y} each time, since they could be integers, floats or lists. The C \omterm{def:b1:compilation-pipeline:compiler}{compiler} made those
decisions once, before the program ran. The price of the \omterm{def:b1:compilation-pipeline:interpreter}{interpreter}'s flexibility can be measured. Here is the
sieve of Eratosthenes, written the same way in both languages.

\omcode{code/bachelor-1/01-compilation-pipeline/listings/sieve.c}{2}{29}{The sieve in C, timed}\label{lst:b1:compilation-pipeline:sieve}

\omcode{code/bachelor-1/01-compilation-pipeline/listings/pysieve.py}{2}{24}{The same sieve in Python}

\omoutput{code/bachelor-1/01-compilation-pipeline/outputs/sieve.txt}{The C sieve, one run on the author's
laptop.}

\omoutput{code/bachelor-1/01-compilation-pipeline/outputs/pysieve.txt}{The Python sieve, one run on the author's
laptop.}

Both print the same count, and the compiled program is more than ten times faster here. That is not a verdict on
the languages. Python started at once, found its errors with readable messages, and its own library offers
faster routes (a slice assignment instead of the inner loop). C asked for a build, would have printed nonsense
for an out-of-range index instead of an error, and gave speed in return. The rest of this volume uses both,
each where it fits: C where the machine is the subject, Python where the algorithm is.

\begin{omfigure}
\begin{tikzpicture}[
  box/.style={draw, rounded corners=2pt, minimum height=8mm, align=center, font=\small},
  >={Stealth[length=2mm]}]
\node[box, fill=omCode!8, minimum width=24mm] (src) at (0,0) {\code{square_plus}\\[-2pt]source};
\node[box, fill=omCode!8, minimum width=24mm] (bc) at (3.8,0) {bytecode};
\node[box, fill=omThm!10, align=left, font=\footnotesize, inner sep=4pt] (loop) at (8.6,0)
  {\textbf{interpreter loop}\\fetch the next instruction\\check the operands' types\\do the operation};
\draw[->] (src) -- node[above, font=\scriptsize] {compile} (bc);
\draw[->] (bc) -- (loop);
\draw[->, thick] ([yshift=3mm]loop.east) arc[start angle=90, end angle=-90, radius=3mm];
\node[box, fill=omCode!8, minimum width=24mm] (csrc) at (0,-2.6) {\code{hello.c}\\[-2pt]source};
\node[box, fill=omCode!8, minimum width=24mm] (exe) at (3.8,-2.6) {machine code};
\node[box, fill=omThm!10, minimum width=34mm] (cpu) at (8.6,-2.6) {processor};
\draw[->] (csrc) -- node[above, font=\scriptsize] {gcc} (exe);
\draw[->] (exe) -- node[above, font=\scriptsize] {executed} (cpu);
\end{tikzpicture}
\par\smallskip{\small Two routes from source to execution. Above, Python compiles to \omterm{def:b1:compilation-pipeline:interpreter}{bytecode} and a loop
interprets it, deciding types at each step. Below, \code{gcc} produces machine code that the processor executes
directly, with every decision taken before the run.}
\end{omfigure}

\section{Exercises}\label{sec:b1:compilation-pipeline:exercises}

\begin{exercise}[$\star$]\label{exo:b1:compilation-pipeline:1}
Which stage of a build, or the running program, reports each problem?
(a) \texttt{\#include "sqare.h"}, misspelt, when only \code{square.h} exists;
(b) a missing semicolon at the end of a statement;
(c) \code{undefined reference to `cube'};
(d) a division by zero whose divisor is read from the keyboard.
\end{exercise}

\begin{exercise}[$\star$]\label{exo:b1:compilation-pipeline:2}
A module \code{geometry.c} was compiled with \code{gcc -c}, and \code{nm} printed its \omterm{def:b1:compilation-pipeline:compiler}{symbols}:
\omoutput{code/bachelor-1/01-compilation-pipeline/outputs/exo_nm.txt}{The \omterm{def:b1:compilation-pipeline:compiler}{symbols} of an unseen \omterm{def:b1:compilation-pipeline:compiler}{object file}.}
For each of the five lines, say whether the name is a function or a variable, defined here or elsewhere, visible to
other files or not, and for the variables whether they start at zero. Then guess the declarations.
\end{exercise}

\begin{exercise}[$\star$]\label{exo:b1:compilation-pipeline:3}
In which region of the running program does each object (a) to (f) of the program below live, and how long?
\omcode{code/bachelor-1/01-compilation-pipeline/listings/where.c}{2}{17}{Six objects to place}
\end{exercise}

\begin{exercise}[$\star$]\label{exo:b1:compilation-pipeline:4}
Write out what the \omterm{def:b1:compilation-pipeline:source-file}{preprocessor} makes of the two \code{printf} lines of the program below, and predict what it
prints.
\omcode{code/bachelor-1/01-compilation-pipeline/listings/macro_trap.c}{2}{12}{A \omterm{def:b1:compilation-pipeline:source-file}{macro} and its careful twin}
\end{exercise}

\begin{exercise}[$\star\star$]\label{exo:b1:compilation-pipeline:5}
The files \code{users.c} and \code{admin.c} both include the header \code{limits_bad.h}, which holds the single
line \code{int max_users = 100;}. Each file compiles, but the link fails:
\omoutput{code/bachelor-1/01-compilation-pipeline/outputs/exo_header.txt}{The build of a two-file program.}
Explain the message, and change the header and one \omterm{def:b1:compilation-pipeline:source-file}{source file} so that the program builds, with one variable
shared by both files.
\end{exercise}

\begin{exercise}[$\star\star$]\label{exo:b1:compilation-pipeline:6}
The program below always finds no negative reading. Read its warning, then fix the program.
\omcode{code/bachelor-1/01-compilation-pipeline/listings/negatives.c}{2}{13}{Counting negative readings}
\omoutput{code/bachelor-1/01-compilation-pipeline/outputs/negatives.txt}{Its build and run.}
\end{exercise}

\begin{exercise}[$\star\star$]\label{exo:b1:compilation-pipeline:7}
A one-file program computes the mean and the maximum of an array of \code{double} with two functions
\code{mean(const double *x, size_t n)} and \code{maximum(const double *x, size_t n)}, and prints them from
\code{main}. Split it into a module \code{stats} (a header with include guards and a \code{.c} file) and a file
\code{stats_main.c}; give the commands that build it in separate steps.
\end{exercise}

\begin{exercise}[$\star\star$]\label{exo:b1:compilation-pipeline:8}
A header \code{point.h} contains only \code{struct point { int x, y; };}, with no include guard. The headers
\code{shapes.h} and \code{draw.h} both include it, and \code{main.c} includes \code{shapes.h} then \code{draw.h}.
How many copies of \code{point.h} does the \omterm{def:b1:compilation-pipeline:source-file}{preprocessor} put into \code{main.c}, what happens next, and what do
include guards change? With $k$ headers each including \code{point.h}, all included by \code{main.c}, how many
copies are there with and without guards?
\end{exercise}

\begin{exercise}[$\star\star$]\label{exo:b1:compilation-pipeline:9}
Predict the three lines printed by this program.
\omcode{code/bachelor-1/01-compilation-pipeline/listings/calls.c}{2}{19}{A static and an automatic local}
\end{exercise}

\begin{exercise}[$\star\star\star$]\label{exo:b1:compilation-pipeline:10}
A program has $k = 40$ \omterm{def:b1:compilation-pipeline:source-file}{source files} of about $L = 500$ lines each, and one header included by $m = 12$ of them.
(a) After an edit to one \omterm{def:b1:compilation-pipeline:source-file}{source file}, how many lines does a rebuild compile with separate \omterm{def:b1:compilation-pipeline:compiler}{compilation}, and how
many if the program were one file? (b) Same question after an edit to the header. (c) A day of work makes 30
edits to \omterm{def:b1:compilation-pipeline:source-file}{source files} and 4 to the header. Count the lines compiled in each organisation, and the ratio.
\end{exercise}

\begin{exercise}[$\star\star\star$]\label{exo:b1:compilation-pipeline:11}
Write a function \code{const char *region_of(const void *p, const void *stack_marker)} that returns
\code{"code"}, \code{"static data"}, \code{"heap"} or \code{"stack"} for a pointer into the corresponding region
of the running program, given the address of any local variable of the caller as \code{stack_marker}. Explain
why such a function can only work on a known layout, and why the C language does not promise it.
\end{exercise}

\begin{exercise}[$\star\star\star$]\label{exo:b1:compilation-pipeline:12}
Three \omterm{def:b1:compilation-pipeline:compiler}{object files} have these \omterm{def:b1:compilation-pipeline:compiler}{symbols}. \code{a.o}: \code{T main}, \code{U parse}, \code{U total}.
\code{b.o}: \code{T parse}, \code{D total}, \code{U emit}. \code{c.o}: \code{T emit}, \code{T parse}.
(a) Which of the links \code{a.o b.o}, \code{a.o b.o c.o}, \code{a.o c.o} succeed, and why do the others fail?
(b) Using \cref{prop:b1:compilation-pipeline:link}, count the subsets of $\{\code{a.o}, \code{b.o},
\code{c.o}\}$ containing \code{a.o} that link successfully.
(c) If \code{parse} were declared \code{static} in \code{c.o}, what would change?
\end{exercise}

\section{Problem: One program, four files, one process}\label{sec:b1:compilation-pipeline:problem}

\begin{problem}[{Weekend problem --- One program, four files, one process}]\label{pb:b1:compilation-pipeline:1}
A small word counter is split into four files: a header \code{count.h}, a module \code{count.c} that counts
lines and words one character at a time, a module \code{report.c} that prints the counts, and \code{wcmain.c},
which reads the standard input.
\omcode{code/bachelor-1/01-compilation-pipeline/listings/count.h}{2}{11}{The shared header}
\omcode{code/bachelor-1/01-compilation-pipeline/listings/count.c}{2}{18}{Counting lines and words}
\omcode{code/bachelor-1/01-compilation-pipeline/listings/report.c}{2}{15}{Printing a report}
\omcode{code/bachelor-1/01-compilation-pipeline/listings/wcmain.c}{2}{18}{The main program}
\omoutput{code/bachelor-1/01-compilation-pipeline/outputs/problem.txt}{Building and running the word
counter.}\label{out:b1:compilation-pipeline:problem}

\needspace{4\baselineskip}
\medskip\noindent\textbf{Part I --- Preprocessing.}
\begin{enumerate}
\item Which files does the \omterm{def:b1:compilation-pipeline:source-file}{preprocessor} read to produce the expanded \code{count.c}? Why is the expansion more
than 300 lines long when \code{count.c} has 18?
\item What does the \omterm{def:b1:compilation-pipeline:source-file}{preprocessor} put in place of \code{MAX_REPORTS} in \code{report.c}? Does the \omterm{def:b1:compilation-pipeline:compiler}{compiler} ever
see the name \code{MAX_REPORTS}?
\item Does \code{count.c} see the declaration of \code{report}? Does that matter?
\item If the include guard of \code{count.h} were removed, would any of the three files fail to compile? Justify
from what the header contains.
\item Someone writes \texttt{\#define MAX\_REPORTS 3;} with a semicolon. Write the line of \code{report.c} that the
\omterm{def:b1:compilation-pipeline:compiler}{compiler} would receive and say whether it still compiles.
\end{enumerate}

\needspace{4\baselineskip}
\medskip\noindent\textbf{Part II --- Object files.}
\begin{enumerate}[resume]
\item List the \omterm{def:b1:compilation-pipeline:compiler}{symbols} that \code{count.o} defines and the one it uses. Where does the used one come from?
\item Why is \code{in_word} printed with a lower-case \code{b} and \code{lines} with an upper-case \code{B}?
\item Why is \code{reports} printed as \code{reports.0}, with a lower-case letter?
\item The variable \code{unit} is a \code{D} \omterm{def:b1:compilation-pipeline:compiler}{symbol} of 8 bytes in \code{report.o}. What are those 8 bytes, and
where are the six bytes of \code{"words"} stored?
\item Explain the \code{bss} column of \code{size}: 20 bytes for \code{count.o}, 4 for \code{report.o}. Why does
\code{in_word}, initialised to 0, count as bss?
\item The file \code{wcmain.o} defines one \omterm{def:b1:compilation-pipeline:compiler}{symbol} and uses six. Classify the six into those another file of the
program defines and those the C library defines.
\end{enumerate}

\needspace{4\baselineskip}
\medskip\noindent\textbf{Part III --- Linking.}
\begin{enumerate}[resume]
\item Check \cref{prop:b1:compilation-pipeline:link} on the three \omterm{def:b1:compilation-pipeline:compiler}{object files}, ignoring the C library: list
each used \omterm{def:b1:compilation-pipeline:compiler}{symbol} and its unique definition.
\item What would the link print without \code{report.o}? Name every \omterm{def:b1:compilation-pipeline:compiler}{undefined reference}.
\item A colleague adds \code{long words;} to \code{wcmain.c} at file level. What happens at the link? And with
\code{static long words;}?
\item Why could the link succeed even though the three \omterm{def:b1:compilation-pipeline:compiler}{object files} were compiled at different times?
\item Could \code{count.c} and \code{report.c} each contain a \code{static int} called \code{in_word}? Two
non-static ones?
\end{enumerate}

\needspace{4\baselineskip}
\medskip\noindent\textbf{Part IV --- The running program.}
\begin{enumerate}[resume]
\item The program printed \code{2 lines, 5 words}. Trace \code{feed} on the input to check the 5.
\item Ten objects take part in a run: the code of \code{feed}, the code of \code{main}, \code{lines},
\code{words}, \code{in_word}, \code{unit}, the characters of \code{"words"}, \code{reports}, the local \code{c} of
\code{main}, and the 16 bytes that \code{title} points to. Place each in its region.
\item For each of the ten, say when its lifetime starts and ends.
\item Which of the ten exist before \code{main} starts, and with which values?
\item If \code{report} were called four times, what would the fourth call print, and why does
\code{reports} remember the earlier calls while a local variable of \code{report} could not?
\item The pointer \code{title} itself is a local variable. In which region is it, and in which region is what it
points to?
\item After \code{free(title)}, which of the ten objects no longer exist, and which still do?
\item Count the objects of the placement table in each region and state the result.
\end{enumerate}
\end{problem}
